% *======================================================================*
%  Cactus Thorn template for ThornGuide documentation
%  Author: Ian Kelley
%  Date: Sun Jun 02, 2002
%  $Header$
%
%  Thorn documentation in the latex file doc/documentation.tex
%  will be included in ThornGuides built with the Cactus make system.
%  The scripts employed by the make system automatically include
%  pages about variables, parameters and scheduling parsed from the
%  relevant thorn CCL files.
%
%  This template contains guidelines which help to assure that your
%  documentation will be correctly added to ThornGuides. More
%  information is available in the Cactus UsersGuide.
%
%  Guidelines:
%   - Do not change anything before the line
%       % START CACTUS THORNGUIDE",
%     except for filling in the title, author, date, etc. fields.
%        - Each of these fields should only be on ONE line.
%        - Author names should be separated with a \\ or a comma.
%   - You can define your own macros, but they must appear after
%     the START CACTUS THORNGUIDE line, and must not redefine standard
%     latex commands.
%   - To avoid name clashes with other thorns, 'labels', 'citations',
%     'references', and 'image' names should conform to the following
%     convention:
%       ARRANGEMENT_THORN_LABEL
%     For example, an image wave.eps in the arrangement CactusWave and
%     thorn WaveToyC should be renamed to CactusWave_WaveToyC_wave.eps
%   - Graphics should only be included using the graphicx package.
%     More specifically, with the "\includegraphics" command.  Do
%     not specify any graphic file extensions in your .tex file. This
%     will allow us to create a PDF version of the ThornGuide
%     via pdflatex.
%   - References should be included with the latex "\bibitem" command.
%   - Use \begin{abstract}...\end{abstract} instead of \abstract{...}
%   - Do not use \appendix, instead include any appendices you need as
%     standard sections.
%   - For the benefit of our Perl scripts, and for future extensions,
%     please use simple latex.
%
% *======================================================================*
%
% Example of including a graphic image:
%    \begin{figure}[ht]
% 	\begin{center}
%    	   \includegraphics[width=6cm]{MyArrangement_MyThorn_MyFigure}
% 	\end{center}
% 	\caption{Illustration of this and that}
% 	\label{MyArrangement_MyThorn_MyLabel}
%    \end{figure}
%
% Example of using a label:
%   \label{MyArrangement_MyThorn_MyLabel}
%
% Example of a citation:
%    \cite{MyArrangement_MyThorn_Author99}
%
% Example of including a reference
%   \bibitem{MyArrangement_MyThorn_Author99}
%   {J. Author, {\em The Title of the Book, Journal, or periodical}, 1 (1999),
%   1--16. {\tt http://www.nowhere.com/}}
%
% *======================================================================*

% If you are using CVS use this line to give version information
% $Header$

\documentclass{article}

% Use the Cactus ThornGuide style file
% (Automatically used from Cactus distribution, if you have a
%  thorn without the Cactus Flesh download this from the Cactus
%  homepage at www.cactuscode.org)
\usepackage{../../../../doc/latex/cactus}

\begin{document}

% The author of the documentation
\author{Erik Schnetter \textless schnetter@gmail.com\textgreater}

% The title of the document (not necessarily the name of the Thorn)
\title{HydroToyAMReX}

% the date your document was last changed, if your document is in CVS,
% please use:
%    \date{$ $Date$ $}
\date{July 10 2019}

\maketitle

% Do not delete next line
% START CACTUS THORNGUIDE

% Add all definitions used in this documentation here
%   \def\mydef etc

% Add an abstract for this thorn's documentation
\begin{abstract}

\end{abstract}

% The following sections are suggestive only.
% Remove them or add your own.

\section{Introduction}

\section{Physical System}

\section{Numerical Implementation}

\section{Using This Thorn}

\subsection{Obtaining This Thorn}

\subsection{Basic Usage}

\subsection{Special Behaviour}

\paragraph{Nodal gauge register (\texttt{AsterX::gauge\_register}).}
In the generalized Lorenz gauge the vector potential receives the gauge term
$-\partial_i G$. Under CarpetX subcycling the coarse and fine levels
accumulate different time integrals of $G$ at the nodes they share, and the
driver's restriction of \texttt{Avec} leaves that mismatch as a spurious curl
(a $B$ shell) in the first coarse cell outside each refinement boundary. With
\texttt{gauge\_register = yes} (the default) AsterX evolves a vertex ledger
\texttt{IG} with $\dot{\texttt{IG}} = G$ on every level and, in
\texttt{CCTK\_POSTRESTRICT} on every time-aligned pass, shifts the coarse
\texttt{Avec} edges near each interface by the edge-difference of
$\texttt{IG}_\mathrm{fine} - \texttt{IG}_\mathrm{coarse}$ (a pure gauge
transformation using the same finite-difference operator as the
$-\partial_i G$ term, at \texttt{mag\_correction\_order}) so both levels use
one gauge integral per shared node. The register never enters the RHS,
changes \texttt{Avec} only within one gauge-stencil width of a refinement
boundary, and is re-zeroed on every full cascade, so \texttt{IG} never holds
more than one coarse step's worth of $\int G\,dt$ on a multi-level run.

\emph{Activation.} The register is active only when all three of
\texttt{gauge\_register = yes}, \texttt{vector\_potential\_gauge =
"generalized Lorenz"} and \texttt{CarpetX::use\_subcycling = yes} hold.
Otherwise $\texttt{IG\_rhs} = 0$, none of the reconcile routines is
scheduled, and the run is bit-identical to one without the register (without
subcycling the mismatch is identically zero). Under subcycling CarpetX
requires \texttt{restrict\_during\_sync = no}, so \texttt{CCTK\_POSTRESTRICT}
is always traversed when the register is active.

\emph{Regridding.} On a two-level run every regrid that modifies a level
happens at a time-aligned point, where the preceding full cascade has already
zeroed \texttt{IG} on every level; the ledger prolongated onto new fine points
is therefore exactly zero and exact, and the correction runs unchanged on the
next full cascade. (With three or more levels a regrid of the finest level at
a partial cascade prolongates a non-zero intermediate ledger; that one-time
interpolation residual is accepted.) On a subcycled run with a single level no
full cascade ever runs, so \texttt{IG} accumulates $\int G\,dt$ over the whole
run without being re-zeroed; this has no effect on \texttt{Avec}, since no
reconcile runs either.

\emph{Cost.} Storage for \texttt{IG}, \texttt{IG\_rhs} and the scratch
\texttt{IGr} is always allocated (ODESolvers requires storage on every
rhs-tagged group), \texttt{IG\_rhs} is written inside the \texttt{Psi} RHS
loop, and \texttt{IG} is synchronized with the other evolved variables after
every stage whether or not the register is active (the driver requires every
evolved group to be valid everywhere at the top of each iteration). An inert
run therefore still pays the ghost synchronization of one vertex grid
function, but no extra kernel and no reconcile.

\emph{Validation and tests.} On a two-level magnetised TOV run with
\texttt{mag\_correction\_order = 2}, the register reduces the interface
$B$-shell metric by a factor of about 2360 ($2.1\times10^{-14}$ to
$9.1\times10^{-18}$) and brings the subcycled interface profile onto the
non-subcycled one. Three regression tests cover the AMR configurations:
\texttt{magTOV\_Z4c\_AMR\_SC} and \texttt{magTOV\_Z4c\_AMR\_SC\_norestrict}
run subcycled with \texttt{gauge\_register = no} and pin the subcycled path
bit-for-bit against reference data that predate the register;
\texttt{magTOV\_Z4c\_AMR\_SC\_gauge\_register} runs the same setup as
\texttt{magTOV\_Z4c\_AMR\_SC} with the register on and CarpetX poisoning at
its default, against reference data generated with the register on;
\texttt{magTOV\_Z4c\_AMR} is not subcycled, so the register is inert there.

\subsection{Interaction With Other Thorns}

\subsection{Examples}

\subsection{Support and Feedback}

\section{History}

\subsection{Thorn Source Code}

\subsection{Thorn Documentation}

\subsection{Acknowledgements}


\begin{thebibliography}{9}

\end{thebibliography}

% Do not delete next line
% END CACTUS THORNGUIDE

\end{document}
